\chapter{Background and Related Work}

\section{Large Language Models}
Large Language Models (LLMs) represent a major advancement in natural language processing. These models, which are generally based on the transformer architecture, are pre-trained on massive collections of text to learn statistical representations of language structure, syntax, and semantics. Through this pre-training, LLMs acquire general capabilities in text generation, question answering, and text summarization, allowing them to process natural language inputs and synthesize coherent responses across various domains.

In domain-specific applications, LLMs are frequently utilized to interpret user queries and generate conversational responses. When applied to question answering, the model leverages its internal parametric memory—the associations and facts encoded within its weights during training—to construct answers. This capability allows models to translate informal, conversational user questions into structured responses or intermediate representations, making them useful interfaces for information access.

However, in legal question answering, relying solely on an LLM's parametric memory introduces severe limitations. The primary challenge is hallucination, where the model generates text that is linguistically plausible but legally incorrect or entirely fabricated. Furthermore, LLMs frequently experience context limits, lack access to updated or specialized domain knowledge, and struggle to produce accurate, verifiable legal citations. Because statutory rules are highly sensitive, a minor error in an article number or coordinate can alter the legal meaning, creating significant compliance risks.

Therefore, legal question-answering systems should not rely only on model memory to formulate answers. An LLM cannot be treated as a legal oracle. This limitation motivates the need for retrieval-based grounding. Rather than forcing the language model to retrieve statutory provisions from its static parameters, the system should locate verified legal passages from an external, structured corpus and present them to the model, ensuring that answer generation is grounded in retrieved legal sources.

\section{Retrieval-Augmented Generation}
Retrieval-Augmented Generation (RAG) is a framework designed to improve the reliability of language models by combining an external document retriever with a generative network \cite{lewis2020retrieval}. The baseline workflow involves an offline indexing stage and an online query stage. During indexing, a target text corpus is divided into small passages, converted into continuous vector representations using an embedding model, and stored in a vector database. At runtime, the system embeds the user's query, retrieves the most semantically similar passages, and assembles them into the model's prompt context to guide the generation layer.

\begin{figure}[htbp]
  \centering
  \resizebox{0.95\textwidth}{!}{\begin{tikzpicture}[
  node distance=8mm and 10mm,
  every node/.style={font=\footnotesize},
  ragbox/.style={block, minimum height=8mm, minimum width=28mm, text width=28mm, align=center},
  ragstore/.style={ragbox, fill=black!10},
  vectorbox/.style={ragbox, fill=blue!8},
  raggroup/.style={container, fill=black!3, inner sep=7pt,
    label={[anchor=west, fill=black!3, inner sep=2pt, font=\footnotesize\bfseries]north west:#1}}
]
  \node[ragbox] (docs) {Documents};
  \node[ragbox, right=of docs] (chunking) {Chunking};
  \node[ragbox, right=of chunking] (embeddings) {Embeddings};
  \node[vectorbox, right=of embeddings] (vector) {Vector Index};

  \node[ragbox, below=19mm of docs] (query) {User Query};
  \node[vectorbox, right=of query] (retriever) {Retriever};
  \node[ragbox, right=of retriever] (context) {Retrieved Context};
  \node[ragbox, right=of context] (generator) {Generator};
  \node[ragbox, right=of generator] (answer) {Answer with Sources};

  \begin{scope}[on background layer]
    \node[raggroup=Offline Indexing, fit=(docs) (chunking) (embeddings) (vector)] {};
    \node[raggroup=Online Question Answering, fit=(query) (retriever) (context) (generator) (answer)] {};
  \end{scope}

  \draw[arrow] (docs) -- (chunking);
  \draw[arrow] (chunking) -- (embeddings);
  \draw[arrow] (embeddings) -- (vector);

  \draw[arrow] (query) -- (retriever);
  \draw[arrow] (retriever) -- (context);
  \draw[arrow] (context) -- (generator);
  \draw[arrow] (generator) -- (answer);

  \draw[arrow] (vector.south) -- (retriever.north);
\end{tikzpicture}
}
  \caption{Standard retrieval-augmented generation pipeline.}
  \label{fig:standard-rag-pipeline}
\end{figure}

The baseline RAG pipeline, illustrated in Figure 2.1, provides several advantages over standard generative answering. By supplying the generator with retrieved source passages, the framework reduces the risk of factual hallucinations, enables source-grounded answers, and provides a clear mechanism for citation tracing. Furthermore, RAG allows system administrators to update the underlying knowledge base dynamically by modifying the vector database, bypassing the expensive and complex training cycles required to retrain LLM parameters.

However, standard flat RAG architectures exhibit major limitations when applied to legal document corpora. In a flat RAG pipeline, documents are split into independent, fixed-size text blocks, which destroys the nested structure of statutory texts. Because flat retrieval treats these chunks as isolated text blocks, it is blind to parent-child hierarchies and cross-document reference networks. For example, standard flat RAG frequently separates a specific clause from its parent article title or document name, causing semantic fragmentation that makes the chunk uninterpretable for both the embedding model and the downstream generator.

Furthermore, relying solely on semantic vector similarity can lead to incomplete or noisy results. Dense embeddings excel at identifying general conceptual synonyms, but they frequently struggle to match exact article coordinates, numerical coordinate tables, or specific cross-referenced implementing decrees. If an essential statutory exception or guiding circular is ranked low by the vector search, it is omitted from the context window, preventing the generator from producing a legally complete answer. This challenge indicates the need for retrieval strategies that actively preserve and navigate legal document structures.

\section{Legal Question Answering and Domain Challenges}
Legal question answering differs fundamentally from open-domain QA because legal answers require strict authority, jurisdiction specificity, temporal validity, and cautious wording \cite{zhong2020legalqa}. A legal assertion is difficult to verify if it is not grounded in the correct legal basis and supported by precise citation traceability. Users need to verify the exact document, article, clause, and point that establishes a right or obligation. A system that summarizes general legal concepts without identifying these statutory anchors cannot be used for compliance or dispute resolution.

A primary challenge in legal QA is the nested, hierarchical structure of legislative texts. In Vietnam, statutory law is strictly organized by document, chapter, section, article, clause, point, and appendix, a structure governed by the Law No. 80/2015/QH13 on Promulgation of Legislative Documents \cite{vietnam2015law80}. To interpret a single legal clause correctly, the system must maintain its structural parent-child hierarchy and align it with its legislative normative rank. Under the principle of hierarchical authority, higher-level primary codes establish broad rights, while secondary administrative decrees and tertiary circulars detail their execution. A lower-level circular must remain strictly compliant with its parent decree and code, and conflicts are resolved in favor of the higher-ranking document.

This structural complexity is clearly illustrated by the motivating case study of this thesis: employment contract termination. The project was initially motivated by the specific rules and disputes surrounding how labor relations end. However, during system development, it was discovered that contract termination questions cannot be resolved in isolation. Resolving termination inquiries naturally requires broad statutory context, including labor contract rules (probation and contract types), employee and employer rights, unlawful termination criteria, severance and compensation calculations, minor worker rules, and female worker protections. Furthermore, retirement age tables in specialized decrees determine when a contract terminates due to retirement, and civil procedure rules in separate codes govern how disputes are litigated in court.

Consequently, the project's scope was expanded from a narrow, topic-specific case study to a broader Vietnamese labor-law QA system over six indexed documents. This transition highlights a key domain challenge: a legal QA system must retrieve not only semantically similar text but the exact, structured network of related legal provisions and guiding regulations. Benchmarks such as LexGLUE \cite{chalkidis2021lexglue} and LegalBench \cite{guha2023legalbench} show the growing importance of evaluating legal language understanding and legal reasoning systems, although this thesis focuses on a smaller Vietnamese labor-law corpus with citation-level retrieval requirements.

\section{Knowledge Graphs and Graph-Augmented RAG}
Knowledge graphs (KGs) represent structured information by modeling entities and their semantic relationships as directed networks. In natural language processing, knowledge graphs are highly effective at representing complex, relational data that vector embeddings alone cannot capture. By storing data as nodes and typed edges, KGs allow systems to execute graph-traversal queries to discover non-obvious paths and structural dependencies across a text corpus.

In this thesis, the knowledge graph is not a stochastic, LLM-generated entity graph; it is a deterministic, document-structure legal knowledge graph. Rather than using generative models to auto-extract loose semantic entities—which introduces significant noise, missing links, and hallucinations—the graph is built programmatically using stable document nesting, regular-expression-based citation patterns, rule-based taxonomy metadata, and official normative ranks. This design choice ensures that the graph remains more reproducible and easier to inspect.

Graph-augmented retrieval (GraphRAG) combines the advantages of vector retrieval and structured knowledge graphs to address the limitations of flat RAG. In this architecture, vector and hybrid searches in Qdrant are utilized to locate initial, high-relevance seed chunks. Starting from these entry points, the retriever executes parameterized Cypher queries in Neo4j to traverse structural parent edges, cross-document reference links, and taxonomic associations. This topological expansion retrieves parent articles, related circulars, implementing decrees, tables, and procedural guidelines that lack direct semantic similarity with the user's initial query, enriching the generator's context window.

This approach exhibits clear distinctions from the GraphRAG framework proposed by Microsoft \cite{edge2024graphrag}. The Microsoft GraphRAG system is designed for corpus-level global summarization, using LLMs to cluster text units into hierarchical communities and generate text summaries for each community. While effective for synthesizing broad narratives, this design is less directly aligned with the citation-level retrieval requirements of this thesis. In contrast, this thesis designs a smaller, deterministic, high-precision legal graph configured for citation-grounded retrieval and local context expansion over labor-law documents, prioritizing exact statutory coordinate matching over global corpus summarization.

\section{Evaluation Methods for RAG Systems}
Evaluating Retrieval-Augmented Generation systems requires measuring multiple system components independently. RAG evaluation frameworks typically separate overall performance into three core dimensions: retrieval quality (assessing the relevance of retrieved context), answer faithfulness or quality (evaluating semantic correctness and naturalness), and citation grounding (checking whether generated claims match source text). Prior RAG evaluation work such as Ragas \cite{es2023ragas} and ARES \cite{saadfalcon2023ares} motivates separating retrieval quality, answer faithfulness, and answer relevance using automated neural evaluators.

In information retrieval, retrieval quality is quantitatively measured using citation-based metadata matching. Standard metrics include Recall@k (measuring the proportion of ground-truth citations recovered), Required Citation Coverage (verifying the complete retrieval of all gold citations), and the Forbidden Citation Violation Rate (measuring the frequency of retrieving obsolete or prohibited documents). Prioritizing high recall and citation coverage is essential in legal QA because missing a mandatory statutory provision prevents downstream generation from producing a complete and well-grounded answer.

For answer and citation evaluation, standard pipelines often rely on human legal expert reviews or LLM-as-Judge frameworks. However, in high-stakes regulatory settings, human expert review remains prohibitively expensive, and using stochastic language models as judges introduces prompt dependencies, high computational costs, and potential evaluation-level hallucinations. These limitations make stochastic evaluation difficult to reproduce or audit during iterative developer testing.

To ensure high reproducibility, transparency, and safety within a controlled environment, this thesis implements a simpler, deterministic, rule-based software validation setup rather than relying on human legal expert reviews or LLM-as-Judge as the final evaluation methods. This deterministic validation engine uses set-theoretic containment and regular expression parsing to programmatically verify that generated citations and article numbers are supported by retrieved contexts. Although this software-level verification is highly reproducible and useful for identifying structural errors, it does not replace legal expert judgment and may miss subtle semantic nuances or complex interpretative disagreements.
